% !TeX root = ../main.tex
\section{Kế hoạch sử dụng vốn theo từng hạng mục và lộ trình giải ngân}

\textbf{1. Bảng phân bổ nguồn vốn 300.000.000 VNĐ (12.000 USD):}
Nguồn vốn gọi được sẽ được cam kết phân bổ tuân thủ nghiêm ngặt theo tỷ lệ ngân sách chi tiết như sau:

\begin{table}[htbp]
    \centering
    \caption{Phân bổ chi tiết nguồn vốn gọi được (Pre-Seed Round)}
    \resizebox{\textwidth}{!}{%
    \begin{tabular}{@{}lp{0.45\textwidth}rr@{}}
        \toprule
        \textbf{Hạng mục đầu tư} & \textbf{Nội dung chi tiết} & \textbf{Số tiền (VNĐ)} & \textbf{Tỷ lệ (\%)} \\
        \midrule
        \textbf{Marketing \& CAC} & Digital Ads, voucher trợ giá đơn đầu, event sinh viên, CLB E-sports. & 90.000.000 & 30,0\% \\
        \textbf{Phát triển Sản phẩm \& Hạ tầng} & Tích hợp Escrow, eKYC, PayOS, kết nối BNPL, Cloud Server \& API. & 80.000.000 & 26,7\% \\
        \textbf{Vận hành \& Nhân sự} & Phụ cấp Team sáng lập, đội ngũ CSKH, kiểm định \& xử lý tranh chấp. & 90.000.000 & 30,0\% \\
        \textbf{Tư vấn Pháp lý} & Hợp đồng điện tử, điều khoản dịch vụ, bản quyền \& tư vấn pháp lý. & 20.000.000 & 6,7\% \\
        \textbf{Quỹ Dự phòng} & Dự phòng nợ xấu, tổn thất tài sản thiết bị \& rủi ro thanh khoản cọc. & 20.000.000 & 6,7\% \\
        \midrule
        \textbf{TỔNG CỘNG} & & \textbf{300.000.000} & \textbf{100,0\%} \\
        \bottomrule
    \end{tabular}%
    }
\end{table}

\textbf{2. Lộ trình giải ngân theo 4 Quý năm đầu:}
Để quản trị rủi ro dòng tiền và tối ưu hiệu quả sử dụng vốn, nguồn vốn 300 triệu VNĐ sẽ được giải ngân theo 4 giai đoạn tương ứng với các cột mốc phát triển của dự án:

\begin{itemize}
    \item \textbf{Quý 1 (Giai đoạn Triển khai MVP \& Khởi động - Giải ngân 40\% $\sim$ 120 triệu VNĐ):}
    Tập trung 60 triệu cho phát triển công nghệ lõi (Escrow, eKYC, PayOS), 30 triệu hoàn thiện pháp lý/hợp đồng điện tử và 30 triệu khởi động chiến dịch Marketing Cold-start.
    
    \item \textbf{Quý 2 (Giai đoạn Thử nghiệm \& Tăng trưởng - Giải ngân 30\% $\sim$ 90 triệu VNĐ):}
    Giải ngân 40 triệu cho Marketing kĩ thuật số (TikTok/Facebook Ads) và trợ giá đơn thuê; 30 triệu duy trì vận hành CSKH; 20 triệu nâng cấp tính năng nền tảng.
    
    \item \textbf{Quý 3 (Giai đoạn Bứt phá Quy mô - Giải ngân 20\% $\sim$ 60 triệu VNĐ):}
    Phân bổ 30 triệu cho các chương trình Marketing sự kiện sinh viên/Gaming Cafe; 20 triệu cho nhân sự vận hành; 10 triệu duy trì hạ tầng máy chủ.
    
    \item \textbf{Quý 4 (Giai đoạn Chuẩn bị Hòa vốn \& Mở rộng B2B - Giải ngân 10\% $\sim$ 30 triệu VNĐ):}
    Dành 15 triệu hoàn thiện tính năng hiển thị ưu tiên/đẩy tin cho đối tác B2B; 15 triệu trích lập bổ sung Quỹ dự phòng tài sản trước khi bước sang Năm 2.
\end{itemize}

\textbf{3. Chỉ số kiểm soát hiệu quả sử dụng vốn (KPIs):}
Hiệu quả giải ngân vốn sẽ được đo lường chặt chẽ hàng tháng thông qua các chỉ số hiệu suất:
\begin{itemize}
    \item \textbf{Chi phí thu hút người dùng thanh toán (CAC):} Giữ mức CAC hiệu dụng dưới 60.000 VNĐ/user thuê thực tế.
    \item \textbf{Tỷ lệ nợ xấu/thất thoát tài sản:} Đảm bảo tỷ lệ tổn thất thiết bị dưới \textbf{1\%} tổng số đơn thuê nhờ cơ chế eKYC + Escrow.
    \item \textbf{Tốc độ quay vòng vốn cọc (Deposit Velocity):} Tỷ lệ chuyển đổi người dùng đăng ký mở ví cọc đạt tối thiểu 30\% tổng số người thuê thực tế.
\end{itemize}
